\section*{Chapitre \ref{ch:b2:liquid-vapour-diagrams} --- Diagrammes binaires liquide--vapeur et distillation}

\begin{solution}{exo:b2:liquid-vapour-diagrams:1}
Benzène \qty{80.1}{\degreeCelsius}, toluène \qty{110.6}{\degreeCelsius}. Le liquide
de composition 0.2 commence à bouillir vers \qty{102}{\degreeCelsius}~; sa première
vapeur a pour composition 0.38.
\end{solution}

\begin{solution}{exo:b2:liquid-vapour-diagrams:2}
$p = 0.5 \times 1.010 + 0.5 \times 0.388 = \qty{0.699}{bar}$~;
$y = 0.505/0.699 = 0.722$.
\end{solution}

\begin{solution}{exo:b2:liquid-vapour-diagrams:3}
$n_V/n = (0.30 - 0.20)/(0.45 - 0.20) = 0.40$~: \qty{4.0}{mol} de vapeur et
\qty{6.0}{mol} de liquide. Vérification~: $6.0 \times 0.20 + 4.0 \times 0.45 = 3.0 = 10
\times 0.30$.
\end{solution}

\begin{solution}{exo:b2:liquid-vapour-diagrams:4}
(a) $v = 2 + 2 - 1 = 3$~: $T$, $p$ et $x$ sont libres. (b) $v = 2 + 2 - 2 = 2$~; à
pression fixée, 1. (c) L'\omterm{def:b2:liquid-vapour-diagrams:azeotrope}{azéotrope} ajoute la relation $x = y$~: $v = 1$~; à pression
fixée, 0 --- l'\omterm{def:b2:liquid-vapour-diagrams:azeotrope}{azéotrope} bout à température fixe, comme un liquide pur,
bien que sa composition change avec la pression.
\end{solution}

\begin{solution}{exo:b2:liquid-vapour-diagrams:5}
Au \omterm{def:b2:liquid-vapour-diagrams:bubble-dew}{point de rosée}, $yp/p_1^* + (1 - y)p/p_2^* = 1$. À \qty{100}{\degreeCelsius}~:
$0.3 \times 1.013/1.80 + 0.7 \times 1.013/0.742 = 1.124 > 1$ (trop froid)~; à
\qty{105}{\degreeCelsius}~: $0.148 + 0.824 = 0.971 < 1$. Par interpolation, $T \approx
100 + 5 \times 0.124/0.153 = \qty{104}{\degreeCelsius}$. La première goutte a pour composition $x =
yp/p_1^* \approx 0.3 \times 1.013/2.0 = 0.15$.
\end{solution}

\begin{solution}{exo:b2:liquid-vapour-diagrams:6}
Les premières gouttes ont la composition de la vapeur au-dessus du liquide en ébullition,
$y = 0.71$, à \qty{92.1}{\degreeCelsius}. Comme le benzène part préférentiellement, le
liquide s'en appauvrit, sa température d'ébullition s'élève et le distillat s'appauvrit
lui aussi~; les dernières gouttes sont presque du toluène pur. La distillation simple est un
seul étage d'équilibre, répété sur un liquide qui ne cesse de changer~: chaque
fraction est un mélange.
\end{solution}

\begin{solution}{exo:b2:liquid-vapour-diagrams:7}
Alimentation à 0.05, du côté de l'eau par rapport à l'\omterm{def:b2:liquid-vapour-diagrams:azeotrope}{azéotrope}~: la tête fournit l'\omterm{def:b2:liquid-vapour-diagrams:azeotrope}{azéotrope}
(0.88 en moles, 95~\% en masse), le bouilleur garde l'eau. Alimentation à 0.95, du côté de
l'éthanol~: la tête fournit encore l'\omterm{def:b2:liquid-vapour-diagrams:azeotrope}{azéotrope}, de plus basse température d'ébullition,
et le bouilleur garde l'éthanol pur.
\end{solution}

\begin{solution}{exo:b2:liquid-vapour-diagrams:8}
Le mélange bout lorsque $0.10 + p^*_{\text{eau}} = 1.013$, c'est-à-dire
$p^*_{\text{eau}} = \qty{0.91}{bar}$, vers \qty{97}{\degreeCelsius}. Rapport massique
$0.10 \times 136/(0.91 \times 18.02) = 0.83$~: \qty{0.83}{g} d'essence par gramme
d'eau.
\end{solution}

\begin{solution}{exo:b2:liquid-vapour-diagrams:9}
À température ambiante, deux phases L$_2$ et L$_1$ (la composition 0.4 est dans la
lacune). Au chauffage, les deux phases bouillent ensemble à la température de
l'\omterm{def:b2:liquid-vapour-diagrams:miscibility-gap}{hétéroazéotrope}, qui reste constante~; la première vapeur a la composition
hétéroazéotropique (le point). Comme 0.4 est du côté de L$_2$ par rapport au point, la phase
L$_1$ disparaît la première~; la température s'élève alors, la phase L$_2$ restante
s'appauvrit en 1 le long de sa \omterm{def:b2:liquid-vapour-diagrams:bubble-dew}{courbe d'ébullition} et la vapeur suit la \omterm{def:b2:liquid-vapour-diagrams:bubble-dew}{courbe de rosée}, jusqu'à ce que
le point atteigne la \omterm{def:b2:liquid-vapour-diagrams:bubble-dew}{courbe de rosée}~: la dernière goutte s'évapore.
\end{solution}

\begin{solution}{exo:b2:liquid-vapour-diagrams:10}
$N_{\min} = \ln(99 \times 99)/\ln 2.47 = 9.19/0.904 = 10.2$, contre 6.5~: passer
de 95~\% à 99~\% de pureté aux deux extrémités coûte plus de la moitié de plateaux
en plus. Chaque facteur supplémentaire sur $x/(1 - x)$ coûte le même nombre de plateaux, et
la pureté se mesure à l'impureté restante, qui ne diminue que géométriquement.
\end{solution}

\begin{solution}{exo:b2:liquid-vapour-diagrams:11}
Le \omterm{def:b2:liquid-vapour-diagrams:binary-diagram}{diagramme isobare} seul ne le donne pas~: il faut $p_1^*$ et $p_2^*$ à
\qty{80}{\degreeCelsius}, tirés des équations d'Antoine (1.010 et
\qty{0.388}{bar}). Alors $p = p_2^* + x(p_1^* - p_2^*)$ et $p = 1/(y/p_1^* + (1 -
y)/p_2^*)$. Le liquide est stable à haute pression (comprimer une vapeur
la condense), et à basse température sur le \omterm{def:b2:liquid-vapour-diagrams:binary-diagram}{diagramme isobare}~: les deux diagrammes
sont inversés l'un par rapport à l'autre.
\end{solution}

\begin{solution}{exo:b2:liquid-vapour-diagrams:12}
Comme $y$ croît avec $x$ et $y(1 - y) > 0$, $dp/dx$ a le signe de $y - x$~:
la pression croît avec $x$ là où la vapeur est plus riche en 1. Pour un \omterm{def:b2:chemical-potential:ideal-mixture}{mélange
idéal}, $y > x$ pour tout $0 < x < 1$ (\cref{prop:b2:liquid-vapour-diagrams:richer})~:
$p$ est strictement monotone, sans extremum, et il n'y a pas d'\omterm{def:b2:liquid-vapour-diagrams:azeotrope}{azéotrope}.
\end{solution}

\begin{solution}{pb:b2:liquid-vapour-diagrams:1}
\textbf{1.} $T = B/(A - \log_{10}1.013) - C$~: benzène $1203.835/4.01253 + 53.226 =
\qty{353.2}{K}$, \qty{80.1}{\degreeCelsius}~; toluène $1343.943/4.07266 + 53.773 =
\qty{383.8}{K}$, \qty{110.6}{\degreeCelsius}.
\textbf{2.} À \qty{363.15}{K}~: $p_1^* = \qty{1.361}{bar}$, $p_2^* = \qty{0.542}{bar}$.
\textbf{3.} $x = (1.013 - 0.542)/(1.361 - 0.542) = 0.575$~; $y = 0.575 \times
1.361/1.013 = 0.773$.
\textbf{4.} $x = (1.013 - 0.742)/(1.800 - 0.742) = 0.256$~; $y = 0.256 \times
1.800/1.013 = 0.455$.
\textbf{5.} \omterm{def:b2:liquid-vapour-diagrams:bubble-dew}{Points de bulle} (0, 110.6), (0.256, 100), (0.575, 90), (1, 80.1)~; \omterm{def:b2:liquid-vapour-diagrams:bubble-dew}{points
de rosée} (0, 110.6), (0.455, 100), (0.773, 90), (1, 80.1) --- le fuseau de la
figure.
\textbf{6.} $\frac12(p_1^* + p_2^*) = 1.013$~: à \qty{90}{\degreeCelsius} 0.952, à
\qty{95}{\degreeCelsius} 1.102~; par interpolation, $90 + 5 \times 0.061/0.150 \approx
\qty{92}{\degreeCelsius}$ (\qty{92.1}{\degreeCelsius} exactement).
\textbf{7.} $v = 2$, soit 1 à pression fixée~: chaque température fixe $x$ et $y$,
c'est pourquoi le diagramme est formé de deux courbes.
\textbf{8.} $x = (1.013 - 0.636)/(1.569 - 0.636) = 0.404$~; $y = 0.404 \times
1.569/1.013 = 0.626$.
\textbf{9.} $n_V = 100 \times (0.500 - 0.404)/(0.626 - 0.404) = \qty{43}{mol}$,
$n_L = \qty{57}{mol}$.
\textbf{10.} Vapeur $43 \times 0.626 = \qty{26.9}{mol}$~; liquide $57 \times 0.404
= \qty{23.0}{mol}$~; total \qty{49.9}{mol} $\approx 50$.
\textbf{11.} $26.9/50 = 54~\%$ du benzène, dans une vapeur à 63~\% seulement.
\textbf{12.} À \qty{100}{\degreeCelsius}, la vapeur en équilibre a pour composition $y =
0.455 < 0.5$~: le point de l'alimentation est au-dessus de la \omterm{def:b2:liquid-vapour-diagrams:bubble-dew}{courbe de rosée}, l'alimentation est entièrement
vaporisée et rien n'est séparé.
\textbf{13.} À son \omterm{def:b2:liquid-vapour-diagrams:bubble-dew}{point de rosée}, \qty{98.8}{\degreeCelsius} (par la méthode de
l'\cref{exo:b2:liquid-vapour-diagrams:5}).
\textbf{14.} $\alpha = p_1^*/p_2^*$~; à \qty{80.1}{\degreeCelsius}, $1.013/0.390 =
2.60$.
\textbf{15.} À \qty{110.6}{\degreeCelsius}, $2.377/1.013 = 2.35$.
\textbf{16.} $\sqrt{2.60 \times 2.35} = 2.47$.
\textbf{17.} $y = xp_1^*/p$ et $1 - y = (1 - x)p_2^*/p$~; diviser.
\textbf{18.} Rien n'est soutiré~: entre deux plateaux, la vapeur qui monte et
le liquide qui descend transportent la même quantité de matière et de chaque constituant,
$V y_k = L x_{k+1}$ avec $V = L$, donc $x_{k+1} = y_k$.
\textbf{19.} Avec $r = x/(1 - x)$~: $r(y_k) = \alpha\,r(x_k)$ et $x_{k+1} = y_k$, donc
$r(x_{k+1}) = \alpha\,r(x_k)$ et $r(x_D) = \alpha^N r(x_B)$~; $N = \ln[r(x_D)/
r(x_B)]/\ln\alpha$.
\textbf{20.} $N_{\min} = \ln(19 \times 19)/\ln 2.47 = 5.889/0.904 = 6.5$.
\textbf{21.} Une colonne doit fournir du produit~: elle fonctionne donc à \omterm{def:b2:liquid-vapour-diagrams:fractional-distillation}{taux de reflux} fini,
ce qui exige davantage de plateaux~; et un plateau réel n'atteint pas l'équilibre (son
efficacité est inférieure à un).
\textbf{22.} Sept marches, la dernière s'achevant à 0.034, au-dessous de 0.05~: 6.5 plateaux en
décompte continu, sept plateaux entiers (le bouilleur comptant pour un).
\textbf{23.} $x = (95/46.07)/(95/46.07 + 5/18.02) = 0.881$.
\textbf{24.} En tête, au mieux l'\omterm{def:b2:liquid-vapour-diagrams:azeotrope}{azéotrope} (0.88 en moles)~; au pied,
l'eau. Cinquante plateaux ne franchissent pas l'\omterm{def:b2:liquid-vapour-diagrams:azeotrope}{azéotrope}.
\textbf{25.} À reflux total, $\boldsymbol{N_{\min} \approx 6.5}$ \omterm{def:b2:liquid-vapour-diagrams:fractional-distillation}{plateaux
théoriques}.
\end{solution}
